%! TEX root = perilakumengobati.tex
% the main TeX file which is intended to compile, :VimtexReload after adjustment
   % this file should be included in the main file
% :h vimtex-tex-root
\documentclass[../manperil.tex]{subfiles}
%ensure the class of docs

\begin{document}

\section{Tekanan Lingkungan dan Perilaku Mengatasi Masalah}

Bagian modul ini membahas hubungan antara tekanan dari lingkungan fisik terhadap manusia serta bagaimana manusia melakukan perilaku penyesuaian atau coping behavior.

\subsection{Pendahuluan: Konsep Dasar Tekanan Lingkungan dan Coping}


Menurut Carver \& Connor-smith, (2010) menjelaskan bahwa StrategiCoping Stres adalah usaha untuk mengatasi atau mengurangi ancaman,bahaya, serta kesulitan yang dihadapi individu. Menuurt Lazarus \& Folkman,(1984) mendefinisikan Strategi Coping Stressebagai upaya kognitif danperilaku yang terus berubah secara dinamis, yang dilakukan individu untukmengelola tuntutan eksternal dan/atau internal yang dinilai membebani ataumelebihi kapasitas sumber daya yang dimiliki oleh individu tersebut.


Definisi Stres Lingkungan yaitu Proses kejadian lingkungan yang mengancam atau menghilangkan kesejahteraan organisme kemudian menimbulkan respon dari organisme tersebut. Definisi Stres: Beban mental yang oleh individu yang bersangkutan ingin dikurangi atau dihilangkan.

Skema Proses Persepsi yaitu Individu mempersepsikan rangsangan dari lingkungan fisik yang akan menentukan apakah kondisi tersebut berada dalam batas optimal atau di luar batas optimal. Kondisi dalam batas toleransi akan menyebabkan individu berada dalam keadaan homeostatis yang sukses. Kondisi di luar ambang toleransi akan menimbulkan stres yang menuntut adanya coping.

Menurut (Santrock, 2007)) menyatakan bahwa Strategi Coping Stres merujuk pada usaha untuk mengelola situasi yang menekan, memperluasupaya dalam memecahkan masalah kehidupan, serta berusaha mengatasi ataumengurangi stres yang dihadapi. Strategi Coping Stresmerupakan suatubentuk usaha individu dalam mengatur atau menyesuaikan diri terhadaptuntutan eksternal maupun internal yang dirasakan menekan atau melampauikapasitas sumber daya yang dimilikinya (Rukmana, 2019) Berdasarkan beberapa definisi yang telah dijelaskan, penelitimenyimpulkan bahwa Strategi Coping Stresmerupakan usaha individu dalammenghadapi dan mengelola stres yang muncul akibat tuntutan eksternalmaupun internal. Proses ini melibatkan respons kognitif dan perilaku yangbertujuan untuk mengurangi atau menghilangkan ketegangan psikologis.Terdapat dua pendekatan utama dalam Strategi Coping Stres, yaitu yangberfokus pada penyelesaian masalah dan yang berfokus pada pengelolaanemosi, yang keduanya bertujuan untuk membantu individu menyesuaikan diridengan kondisi yang penuh tekanan


\subsection{Dimensi Strategi Coping Stres}%
\label{sub:Dimensi Strategi Coping Stres}

Menurut Carver et al., (1989) merancang alat ukur COPE untukmengidentifikasi berbagai Strategi individu dalam menghadapi stres.Instrumen ini mengelompokkan coping ke dalam tiga dimensi utama, yaitu problem-focused coping, emotion-focused coping, dan less useful coping,yang masing-masing terdiri dari sejumlah subdimensi spesifik, yakni:

\subsubsection{Problem-focused coping}%
\label{ssub:problem-focused_coping}

Problem-focused coping adalah Strategi untuk mengatasi stresdengan fokus pada pemecahan masalah, dilakukan ketika individumerasa dapat mengubah situasi atau memiliki sumber daya untukmenghadapinya. Strategi ini mencakup beberapa subdimensi, yaitu:

\begin{enumerate}
	\item Active CopingIndividu secara langsung mengambil langkah nyata untukmenghadapi atau mengurangi stresor. Ini meliputi upayameningkatkan usaha, melakukan tindakan bertahap, sertamengarahkan energi untuk menyelesaikan masalah secarakonstruktif.
	\item PlanningIndividu merancang langkah-langkah yang perlu dilakukanuntuk menghadapi situasi stres. Proses ini mencakup pemikiranStrategis, merencanakan tahapan penyelesaian masalah, danmemilih metode terbaik untuk menghadapi situasi tersebut.
	\item Seeking Instrumental Social SupportIndividu mencari bantuan yang bersifat praktis atau informatifdari lingkungan sosial. Hal ini termasuk mencari saran, informasi,atau bantuan langsung dari orang lain yang dapat membantumenyelesaikan masalah yang dihadapi.
\end{enumerate}

\subsubsection{Emotion-Focused Coping}%
\label{ssub:emotion-focused_coping}

Emotion-Focused Coping merupakan Strategi koping yangbertujuan untuk mengelola tekanan emosional yang timbul akibatsituasi stres. Strategi ini digunakan ketika individu merasa tidak mampumengubah situasi yang menekan dan lebih memilih untukmenyesuaikan diri secara emosional. Fokus utamanya adalahmengurangi perasaan negatif daripada menyelesaikan masalah secara langsung. Emotion-focused coping terdiri dari beberapa subdimensi,yaitu:

\begin{enumerate}
	\item Seeking Social Support for Emotional ReasonsIndividu mencari dukungan dari orang lain untuk mendapatkanpenguatan secara emosional, seperti simpati, pengertian, ataukenyamanan. Strategi ini membantu individu merasa tidak sendiriandalam menghadapi stres.
	\item Positive ReframingIndividu mencoba melihat situasi dari sisi yang lebih positif, ataumengambil hikmah dari peristiwa yang terjadi. Strategi ini dapatmengarahkan individu untuk lebih mudah melakukan coping yangberorientasi pada penyelesaian masalah.
	\item DenialIndividu menolak mengakui keberadaan stresor atau bertindakseolah-olah masalah tersebut tidak ada. Meskipun dapat mengurangistres sementara, Strategi ini berisiko memperburuk situasi karenamasalah tidak segera ditangani.
	\item AcceptanceIndividu secara sadar menerima kenyataan dari situasi stres yangdihadapi, tanpa mencoba menyangkal atau melarikan diri.Penerimaan ini memungkinkan individu untuk menyesuaikan dirisecara emosional dan bertindak lebih adaptif.
	\item Turning to ReligionIndividu menyerahkan persoalan yang dihadapi kepada Tuhanatau menjalani pendekatan spiritual sebagai bentuk penguatan diri.Strategi ini dianggap penting karena agama sering kali menyediakanrasa nyaman, harapan, dan Social Support yang bermakna bagibanyak orang.
\end{enumerate}
\subsubsection{Less Useful Coping}%
\label{ssub:less_useful_coping}
Less Useful Coping adalah Strategi koping yang dianggap kurangadaptif karena cenderung mengurangi upaya aktif dalam menyelesaikan masalah dan mengalihkan perhatian individu ke aktivitas lain yangtidak membantu menyelesaikan stresor secara langsung. Strategi iniberpotensi menghambat penggunaan koping yang lebih konstruktifseperti active coping. Aspek ini terdiri dari beberapa dimensi, yaitu:

\begin{enumerate}
	\item VentingMengekspresikan emosi negatif seperti kemarahan atau frustasisecara berlebihan. Meskipun memberi ruang untuk mengeluarkanperasaan, fokus yang berlebihan pada emosi dapat menghambatpemecahan masalah secara efektif.
	\item BehavioralIndividu secara bertahap mengurangi atau bahkan menghentikanusaha untuk menghadapi stresor. Hal ini bisa tercermin dalam sikapputus asa atau merasa tidak berdaya, yang pada akhirnyamenyebabkan individu menyerah sebelum mencoba mengatasimasalah.
	\item Self-DistractionIndividu mengalihkan perhatian dari stresor dengan melakukankegiatan lain yang tidak berhubungan, seperti menonton televisi,tidur berlebihan, atau melamun. Strategi ini bersifat sementara dantidak menyelesaikan akar permasalahan.
	\item Substance UseIndividu mengandalkan konsumsi zat adiktif (seperti alkohol atauobat-obatan) untuk meredakan ketidaknyamanan emosional akibatstres. Meskipun memberikan rasa nyaman sesaat, Strategi ini dapatmenimbulkan dampak negatif jangka panjang.
	\item HumorIndividu merespons stresor dengan cara menjadikannya bahanlelucon atau mengejek situasi yang menekan. Walaupun dapatmeredakan ketegangan sesaat, humor yang tidak disertaipenyelesaian masalah dapat menghindarkan individu darimenghadapi situasi secara realistis.
	\item Self-Blame Individu memandang dirinya sebagai penyebab utama darimasalah yang terjadi. Sikap ini bisa memunculkan rasa bersalahyang berlebihan dan menurunkan rasa percaya diri dalammenghadapi tantangan.
\end{enumerate}

Dari uraian di atas, dapat disimpulkan bahwa Strategi Coping Stresmenurut Carver et al. (1989) terdiri dari tiga dimensi utama, yaitu problem-focused coping, emotion-focused coping, dan less useful coping. Ketiganyamencerminkan cara individu menghadapi stres, baik dengan menyelesaikanmasalah, mengelola emosi, maupun merespons secara kurang adaptif.Pemahaman terhadap ketiga Strategi ini penting untuk mendorongpenggunaan coping yang lebih efektif.




\subsection{Komponen Stres Lingkungan}
\begin{enumerate}
	\item Stres terdiri atas 3 komponen yang saling terkait, yaitu stressor, proses, dan respon stres.
	\item Stressor: Pemicu stres lingkungan atau sumber stimulus yang mengancam kesejahteraan seseorang.
	\item Respon stres: Reaksi yang timbul akibat stressor yang melibatkan komponen emosional, pikiran, fisiologis, dan perilaku.
	\item Proses: Transaksi antara stressor dengan kapasitas diri individu.
	\item Stres negatif akan muncul jika sumber stres lebih besar dari kapasitas diri.
	\item Stres positif akan muncul jika sumber tekanan sama atau lebih kecil dari kapasitas diri.
\end{enumerate}
\subsection{Identifikasi Sumber Stres Lingkungan}
\begin{enumerate}
	\item Bencana alam dan bencana teknologi.
	\item Kebisingan dan aktivitas pulang pergi kerja.
	\item Termal atau suhu lingkungan.
	\item Density atau kepadatan, yaitu banyaknya jumlah manusia dalam suatu ruang tertentu.
	\item Crowding, yaitu respon objektif terhadap ruang yang sesak.
	\item Pencahayaan, sirkulasi, dan dimensi ruang.
\end{enumerate}

\subsection{Mekanisme Coping dan Hasil Adaptasi}
\begin{enumerate}
	\item Jika coping yang dilakukan individu berhasil, maka akan terjadi adaptasi atau penyesuaian yang membawa pada efek lanjutan yang sukses.
	\item Jika coping gagal, stres akan berlanjut dan dapat menimbulkan efek lanjutan negatif seperti gangguan jiwa atau penyakit.
\end{enumerate}

\subsection{Faktor yang memengaruhi Strategi Coping Stres}%
\label{sub:faktor_yang_memengaruhi_strategi_coping_stres}

Menurut Smet (2020) terdapat sejumlah faktor yang memengaruhibagaimana individu memilih dan menggunakan Strategi Coping Stres.Faktor-faktor tersebut meliputi usia, tingkat pendidikan, status sosialekonomi, Social Support yang diterima, serta karakteristik kepribadianseperti ketabahan (hardiness), keyakinan terhadap kendali pribadi atasperistiwa (locus of control), tingkat daya tahan (resiliensi), dan pengalamanhidup sebelumnya. Dalam penelitian lainya menyatakan bahwa Faktor-faktoryang memengaruhi Strategi Coping Stresmeliputi kondisi kesehatan danenergi fisik, keyakinan positif, kemampuan dalam memecahkan masalah,keterampilan sosial, Social Support , serta ketersediaan sumber dayamateriil.(Nuriza et al., 2023).

Menurut Mu’tadin (dalam Astuti, 2023), Strategi Coping Stres individudipengaruhi oleh berbagai sumber daya pribadi yang berperan penting dalammenghadapi tekanan. Faktor-faktor tersebut meliputi:

\begin{enumerate}
	\item Kesehatan fisik,
	\item[] Berfungsi sebagai modal dasar karena menghadapi stresmembutuhkan energi dan kondisi tubuh yang prima.
	\item Keyakinan positif
	\item[] Mencerminkan pandangan individu terhadap situasi stres. Individudengan keyakinan positif, terutama dengan internal locus of control,cenderung lebih mampu mengelola stres secara efektif dibandingkanmereka yang merasa tidak berdaya (helplessness) akibat external locus ofcontrol.
	\item Keterampilan memecahkan masalah,
	\item[] Kemampuan untuk mengidentifikasi masalah, mengumpulkaninformasi, mengevaluasi alternatif solusi, dan mengambil keputusan yangtepat untuk mengatasi situasi stres.
	\item Keterampilan sosial,
	\item[] Mencakup kemampuan individu dalam berinteraksi secara efektifsesuai norma sosial, sehingga dapat menjalin hubungan yang mendukung.
	\item Social support,
	\item[] Berupa bantuan emosional dan informasi dari orang tua, keluarga,teman, dan lingkungan, yang dapat memperkuat kapasitas individu dalammenghadapi tekanan.
	\item[] Sumber daya materiil,Dukungan berupa uang, barang, atau layanan yang memungkinkanindividu memenuhi kebutuhan dasar dalam situasi stres.
\end{enumerate}
Berdasarkan berbagai pandangan ahli, dapat disimpulkan bahwaStrategi Coping Stres individu dipengaruhi oleh kombinasi faktor personal,sosial, dan lingkungan. Faktor personal meliputi usia, kesehatan fisik, tingkatpendidikan, keyakinan positif, locus of control, resiliensi, serta keterampilanmemecahkan masalah. Faktor sosial mencakup Social Support (SocialSupport) dari orang tua, keluarga, teman, maupun lingkungan sekitar, sertaketerampilan sosial dalam membangun relasi yang sehat. Sementara itu,faktor lingkungan dan sumber daya eksternal, seperti status sosial ekonomidan ketersediaan sumber daya materiil, juga berperan penting dalammenentukan bagaimana individu memilih dan menggunakan Strategi coping.Dengan demikian, coping stres merupakan hasil interaksi antara kondisi internal individu dan dukungan eksternal yang tersedia, yang bersama-samamemengaruhi efektivitas individu dalam menghadapi tekanan hidup.




\bibliographystyle{apalike}
\bibliography{../filename.bib} % required for citecompletion, not work in mainfile
\end{document}

